\section{Frobenius and comparison with finite coefficients}\label{app:second-frobenius}
This appendix treats the characteristic-primary part over fields $F$ of characteristic $p$ with $[F:F^p]=p$, which include local and global function fields; such a field $F$ is fixed throughout.
Its main result, Proposition~\ref{prop:app-frobenius}, compares twisted and ordinary cycle complexes in weights $w\ge1$ when $D$ has $p$-power period, and shows that $K_i(D;\Z/p^a)=0$ for $i\ge2$ for every central division algebra $D$; this is the input used at the characteristic prime in Sections~\ref{sec:paper-local-norm}, \ref{sec:paper-global-weight-one}, \ref{sec:global-chow}, and~\ref{sec:global-spectrum}.

The comparison of cycle complexes uses the coordinate Frobenius on algebraic simplices, whose pullback produces multiplicities divisible by the index.
The $K$-theoretic consequence also requires the support tower over an imperfect field, which we establish in Lemma~\ref{lem:app-twisted-tower}.
For a central division algebra $D$ over $F$ and $w\ge0$, by the reduced rank (Definition~\ref{def:paper-twisted-chow}), the reduced cycle map $\rho$ identifies $z^w(F,\bullet;D)$ with the subcomplex of $z^w(F,\bullet)$ consisting of the cycles whose coefficient at each generator $[W]$ is divisible by $\ind(D_{\kappa(W)})$.

\begin{lemma}\label{lem:app-frobenius-homotopy}
Let $D$ be a central division algebra over $F$ whose period $q=p^c$ is a power of $p$, and for $n\ge0$ let $f_n:\Delta^n\to\Delta^n$ be the finite flat map that raises every barycentric coordinate to the $q$th power.
In each weight $w\ge1$, the pullbacks $f_n^*$ send the ordinary cycle complex $z^w(F,\bullet)$ into the subcomplex $\rho(z^w(F,\bullet;D))$ of cycles whose coefficient at each generator $[W]$ is divisible by $\ind(D_{\kappa(W)})$.
The identity of the ordinary cycle complex is simplicially homotopic to $f_\bullet^*$ through maps preserving this subcomplex.
\end{lemma}
\begin{proof}
Let $W$ be a generator of codimension $w$, let $L$ be the reduced function field of the inverse image of $W$, and let $e$ be the multiplicity of this inverse image.
The map $f_n$ is finite, flat, and radicial of degree $q^n$, and the function field of $W$ is identified with $FL^q$; hence $e[L:FL^q]=q^n$.
Since $L$ has transcendence degree $n-w$ over $F$ and $[F:F^p]=p$, invariance of the $p$-rank under finite extensions gives $[L:L^q]=q^{1+n-w}$.
The $q$th-power isomorphism carries $LF^{1/q}/L$ to $FL^q/L^q$, so $[L:FL^q]=[L:L^q]/[LF^{1/q}:L]$.
Combining these equalities, we obtain
\[
 e=q^{w-1}[LF^{1/q}:L].
\]
The last degree is a multiple of $\ind(D_L)$: under the $q$th-power isomorphism $F^{1/q}\to F$, restriction acts on Brauer classes by multiplication by $q$, since it raises the defining $2$-cocycles to the $q$th power, so $LF^{1/q}$ splits $D_L$.
Since $w\ge1$, we obtain $\ind(D_L)\mid e$, so $f_n^*$ sends the ordinary cycle complex into the subcomplex with index-divisible coefficients.

For a monotone map $\alpha:[n]\to[1]$, put $s=\sum_{\alpha(i)=1}t_i$ and $h_\alpha(t)_j=(1-s)t_j+st_j^q$.
The sum of the image coordinates is $(1-s)+s\sum_jt_j^q=(1-s)+s(\sum_jt_j)^q=1$.
On a geometric fiber, the equation $\sum_{\alpha(i)=1}h_\alpha(t)_i=s-s^2+s^{q+1}$ leaves only finitely many possibilities for $s$, and for each fixed $s$ only finitely many for each coordinate.
Thus $h_\alpha$ and, by the same calculation, its restrictions to every face are quasi-finite maps between smooth simplices of the same dimension; since the sources are Cohen--Macaulay and the targets are regular, they are flat by miracle flatness \cite[Tag 00R4]{StacksProject}.
Flatness on each face and the composition formula for Tor lengths of properly intersecting cycles along Cartier faces give the compatibility of $h_\alpha^*$ with the face maps \cite[\S\S5.6, 5.9]{KahnLevine2010}.
The maps commute with the sums of coordinates occurring in degeneracies because the $q$th-power map is additive.
Thus $h_\alpha^*$ induces a simplicial homotopy between the identity and $f_\bullet^*$.
Since the index over every residue field after flat pullback divides the original index, $h_\alpha^*$ preserves the subcomplex with index-divisible coefficients.
\end{proof}

\begin{lemma}\label{lem:app-twisted-tower}
Let $k$ be an infinite field, and let $A$ be a central simple algebra over $k$.
For every $a\ge0$, the $a$th layer of the homotopy coniveau tower of $K(-;A)$ at $\Spec k$ is identified, through the cycle map of \cite[\S6.2]{KahnLevine2010}, with the spectrum associated with the twisted cycle complex of $A$ in weight $a$.
In particular, the strongly convergent spectral sequence of \cite[Remark 6.1.2, Corollary 6.1.4]{KahnLevine2010} is available over $k$ without assuming that $k$ is perfect.
\end{lemma}
\begin{proof}
Perfectness enters the cited identification only through the geometric purity of \cite[Corollaries 4.2.3--4.2.4]{LevineHomotopyConiveau} and through \cite[\S5.2]{LevineHomotopyConiveau}; we replace the first by localization in twisted $G$-theory and use only the arguments of the second.
The other inputs, namely \cite[Corollary 3.2.2, Theorem 4.1.1]{LevineHomotopyConiveau}, the moving lemma \cite[Theorem 3.3.5]{LevineMovingLemma}, and \cite[Corollary 6.3.3, Proposition B.5]{KahnLevine2010} together with the computation in the proof of \cite[Lemma 6.4.3]{KahnLevine2010}, require only an infinite base field.
Every residue field of a $k$-scheme contains $k$ and is therefore infinite, and the towers below over a residue field $\kappa(z)$ are taken with $\kappa(z)$ as base field.

For a $k$-scheme $X$ and $d\in\Z$, let $G_{(d)}(X;A)$ be the realization of the simplicial spectrum whose $n$-simplices are the colimit, over closed subsets $W\subset X\times\Delta^n$ with $\dim(W\cap(X\times\Delta^f))\le d+\dim\Delta^f$ for every face $\Delta^f$, of the $G$-theory of $A$-modules supported on $W$.
For $X$ of pure dimension $m$, write $G_a(X)=G_{(m-a)}(X;A)$ and $\mathrm{gr}_a=\cofib(G_{a+1}\to G_a)$.
For $b\ge0$, the \emph{$b$th outer stage} $G_a^{(b)}$ of $G_a$ is the realization of the simplicial spectrum whose $n$-simplices are the colimit, over closed subsets $W\subset X\times\Delta^n$ with $\operatorname{codim}_{X\times\Delta^f}(W\cap(X\times\Delta^f))\ge b$ for every face $\Delta^f$, of $\fib\bigl(G_a(X\times\Delta^n)\to G_a((X\times\Delta^n)\setminus W)\bigr)$; the support conditions inside $G_a$ and the outer support $W$ are imposed separately.
More generally, for a functor $E$ to which this outer construction applies, we write $E^{(b)}$ for its $b$th outer stage and $E^{(b/b+1)}=\cofib(E^{(b+1)}\to E^{(b)})$ for its \emph{$b$th outer layer}; in particular $E^{(0/1)}=\cofib(E^{(1)}\to E^{(0)})$.
We put $Q_{a,b}=G_a^{(b/b+1)}$.
For a closed subscheme $Z\subset X$ of pure codimension $c$ in $X$, \cite[Corollary 6.3.3]{KahnLevine2010} gives $(G_a)_Z(X)\simeq G_{a-c}(Z)$.
Put $Z_n=Z\times\Delta^n$; taking the fiber of restriction twice, the outer support becomes $W\cap Z_n$, which has codimension at least $b-c$ on each face of $Z_n$, and conversely every such support on $Z_n$ arises from some $W$.
From now on, let $X$ be smooth of pure dimension over $k$; the closed subscheme $Z$ need not be smooth.
Admissible supports on $U=X\setminus Z$ need not have admissible closures in $X$; we first compute the double fiber with the supports inherited from $X$ and then, after realization, compare with all admissible supports on $U$ by \cite[Corollary 3.2.2]{LevineHomotopyConiveau} applied to $E=G_a$.
Then $G_a$ is a coniveau stage of $K(-;A)$ and the excision of \cite[Theorem 4.1.1]{LevineHomotopyConiveau} applies.
Applying the same localization to each outer support $W\subset Z\times\Delta^n$ shifts the codimension conditions on every face from $(a,b)$ to $(a-c,b-c)$; closed pushforward and pullback to faces are compatible because the product squares obtained by pulling back $Z\times\Delta^n$ to the faces of the simplex are Tor-independent.
Filtering by codimension, measured in $X$, and passing to generic points, the supported term of the $b$th outer layer at a point $z$ of codimension $c$ is the term with indices $(a-c,b-c)$ over $\kappa(z)$.
To pass to generic points, we remove closed subsets of higher codimension from the ambient smooth $X$, take differences of closed supports, and pass to the limit over open subsets; this does not require the closed strata to be smooth.
This replaces the geometric purity used in \cite[Corollaries 4.2.3--4.2.4]{LevineHomotopyConiveau}.

We claim that the outer layer $Q_{a,b}$ vanishes for $0\le b<a$, by induction on $a$.
At a point $z$ of codimension $c>0$, the supported term of $Q_{a,b}$ is $Q_{a-c,b-c}$ over $\kappa(z)$.
It vanishes if $b-c<0$, since the support conditions of the two adjacent outer stages then coincide, and by induction if $b-c\ge0$, since then $0\le b-c<a-c<a$.
Shrinking the dimension conditions on finitely many faces to a dense open subset justifies the passage to generic points.
For $b=0$, the semilocal-simplex computation in the proof of \cite[Lemma 5.2.1]{LevineHomotopyConiveau}, which uses supports avoiding the vertices and the moving lemma, applied to the outer $0$th layer $G_a^{(0/1)}$ with both the inner and outer face conditions imposed, gives $Q_{a,0}=0$ directly.
For $b>0$, the induction shows that $Q=Q_{a,b}$ is birational, so its supported fibers in positive codimension vanish, and $\mathbb A^1$-invariance gives $Q(K)\simeq Q^{(0/1)}(K)$.
By exactness of taking layers, the right-hand side is the cofiber of $(G_a^{(b+1)})^{(0/1)}\to(G_a^{(b)})^{(0/1)}$, and both terms vanish by the same semilocal-simplex computation, applied to $G_a^{(b)}$ and $G_a^{(b+1)}$ with both the inner and outer face conditions imposed.
This is the argument of \cite[Proposition 5.2.2]{LevineHomotopyConiveau}, with the geometric Gersten terms replaced by the supported terms above.
Separately, the supported term of $\mathrm{gr}_a$ at a point of codimension $c>a$ vanishes, since $G_{a-c}$ and $G_{a+1-c}$ are then both $G$-theory and the dimension conditions agree.

For smooth $X$, applying the outer $a$th layer to $G_a\to K(-;A)$ gives an equivalence on supports of codimension at least $a$.
Applying it to $G_a\to\mathrm{gr}_a$ also gives an equivalence, because the contribution of $G_{a+1}$ vanishes by the claim; the outer layers with $b<a$ vanish as well, since $Q_{a,b}=Q_{a+1,b}=0$.
Since $\mathrm{gr}_a$ vanishes in outer codimension at least $a+1$, the outer $a$th layer recovers $\mathrm{gr}_a$.
To compute it, note that in each simplicial degree $n$, the vanishing in higher codimension removes intersections of supports, and additivity for finitely many disjoint supports gives the direct sum of $\mathrm{gr}_0(\kappa(w);A_{\kappa(w)})$ over the points $w$ of codimension $a$ in $X\times\Delta^n$ whose closures satisfy all the face conditions.
Each term is identified with $K(\widehat\Delta^\bullet;A_{\kappa(w)})$, the $K$-theory of the semilocal simplices over $\kappa(w)$, localized for supports avoiding the vertices.
The semilocal-simplex computation in the proof of \cite[Lemma 6.4.3]{KahnLevine2010} gives $\mathrm{gr}_0(\kappa(w);A_{\kappa(w)})\simeq H\bigl(K_0(A_{\kappa(w)})\bigr)$: the other degrees vanish by the computation with the full semilocal simplex and its boundary, and $\pi_0=K_0(A_{\kappa(w)})$ because the two vertex evaluations of a semilocal $1$-simplex agree in reduced rank.
This computation uses only the $K_1$-regularity of the boundary from Vorst's theorem and \cite[Proposition B.5]{KahnLevine2010}, surjectivity of $\operatorname{GL}$ from the semilocal ring to the boundary, and injectivity on $K_0$, all of which hold over an arbitrary residue field.
We use only this computation of $HK_0$, and not the sheaf identification in that lemma, which relies on \cite[Theorem 3.2.4]{KahnLevine2010}.
The map on $\pi_0$ is restriction to the generic point, which is the cycle map of \cite[\S6.2]{KahnLevine2010}.
In summary, for smooth $X$ we have
\[
 K(-;A)^{(a/a+1)}(X)\xleftarrow{\ \sim\ }G_a^{(a/a+1)}(X)\xrightarrow{\ \sim\ }\mathrm{gr}_a^{(a/a+1)}(X)\xleftarrow{\ \sim\ }\mathrm{gr}_a(X).
\]
\end{proof}

\begin{proposition}\label{prop:app-frobenius}
Let $D$ be a central division algebra over $F$.
\begin{enumerate}
\item If the period of $D$ is a power of $p$, then in each weight $w\ge1$ the reduced cycle map from the twisted cycle complex of $D$ to the ordinary cycle complex of $F$ is a quasi-isomorphism with integral coefficients.
\item For $i\ge2$ and every $a\ge1$, one has $K_i(D;\Z/p^a)=0$, and multiplication by $p$ is an isomorphism on $K_i(D)$.
\end{enumerate}
\end{proposition}
\begin{proof}
(1) As recalled before Lemma~\ref{lem:app-frobenius-homotopy}, $\rho$ identifies the twisted cycle complex with the subcomplex $\rho(z^w(F,\bullet;D))$ of $z^w(F,\bullet)$; the identification is compatible with the face maps by \cite[Lemma 5.9.1]{KahnLevine2010}.
By Lemma~\ref{lem:app-frobenius-homotopy}, the identity of the ordinary cycle complex is simplicially homotopic to $f_\bullet^*$ through maps preserving this subcomplex, so the homotopy descends to the quotient, on which $f_\bullet^*$ induces zero.
On the alternating face complex, it gives a chain homotopy from the identity of the quotient to zero, so the quotient is acyclic.

(2) Since $[F:F^p]=p>1$, the field $F$ is imperfect and hence infinite.
The prime-to-$p$ Brauer components can be removed by restriction to a finite separable extension of degree prime to $p$ that splits them; this preserves $[F:F^p]=p$, and for comparison modulo $p$ the restriction is a split injection, since the composite of restriction and corestriction is multiplication by the degree.
By Morita equivalence we may then assume that $D$ has $p$-power period.
By Lemma~\ref{lem:app-twisted-tower}, the strongly convergent spectral sequence for the twisted cycle complexes of $D$ is available over $F$, and the connectivity estimates for the tower remain valid with Moore coefficients.
In weight $w\ge1$, part~(1) identifies the mod-$p$ layers with the commutative ones, which are concentrated in degree $w$, where they are $K_w^M(F)/p$ \cite{GeisserLevine}.
For $w\ge2$, they vanish: $K_w^M(F)/p$ embeds into $\Omega^w_F$ by the $d\log$ map \cite[Corollary (2.8)]{BlochKato1986}, and $\Omega^w_F=0$ since $[F:F^p]=p$.
The weight-zero layer is concentrated in degree zero.
Hence $K_i(D;\Z/p)=0$ for $i\ge2$.
The coefficient sequences for $\Z/p\to\Z/p^{a+1}\to\Z/p^a$ give $K_i(D;\Z/p^a)=0$ for $i\ge2$ and all $a\ge1$ by induction on $a$.
In the Bockstein sequence $0\to K_i(D)/p\to K_i(D;\Z/p)\to K_{i-1}(D)[p]\to0$, the vanishing for $i\ge2$ gives $K_i(D)/p=0$ for $i\ge2$ and $K_j(D)[p]=0$ for $j\ge1$.
Thus multiplication by $p$ is an isomorphism on $K_i(D)$ for $i\ge2$.
\end{proof}
